\section{Rekomendasi Tindak Lanjut}

\subsection{Studi Lanjutan}

\begin{itemize}
  \item \textbf{Kriptografi Lanjutan}: ECC, post-quantum cryptography (lattice-based, hash-based)
  \item \textbf{Keamanan Jaringan}: Implementasi TLS, serangan dan mitigasi
  \item \textbf{Keamanan Aplikasi}: Secure coding, vulnerability assessment
\end{itemize}

\subsection{Sumber Belajar}

\begin{itemize}
  \item \textbf{Buku}: Katz \& Lindell, Handbook of Applied Cryptography, Stallings
  \item \textbf{Standar}: NIST Cryptographic Standards, IETF RFCs
  \item \textbf{Online}: The Joy of Cryptography, OpenSSL, NIST CSRC
  \item \textbf{Praktik}: Proyek enkripsi file, CTF (Capture The Flag)
\end{itemize}

\subsection{Tindak Lanjut}

Terus praktik implementasi dalam C; gunakan OpenSSL untuk RSA, AES, SHA. Ikuti perkembangan standar kriptografi (NIST, IETF). Pertimbangkan sertifikasi (CompTIA Security+, CISSP, CEH) dan bergabung dengan komunitas keamanan informasi \cite{stallings2017}.
